\documentclass{article}
\usepackage[UTF8]{ctex}
\usepackage{amsmath,amssymb,amsfonts}
\usepackage{enumitem}
\usepackage{geometry}
\geometry{a4paper, margin=1in}
% === 错题本排版优化（不参与题目标识） ===
\setlength{\parindent}{0pt}
\setlength{\parskip}{0.35em plus 0.10em minus 0.08em}
\setlength{\emergencystretch}{3em}
\setlist{itemsep=0.25em, topsep=0.40em, parsep=0pt}
\allowdisplaybreaks[2]
% === 错题本排版优化结束 ===
\begin{document}
\section*{东华大学2025年高等代数真题}
\section*{42．东华大学2025年高等代数试卷（回忆版）}

\subsection{题目1}
1．计算行列式

\[
D=\left|\begin{array}{cccc}
a_{1} & x & \cdots & x \\
x & a_{2} & \cdots & x \\
\vdots & \vdots & \ddots & \vdots \\
x & x & \cdots & a_{n}
\end{array}\right| .
\]

\subsection{题目2}
2．设 \(A\) 为 \(m \times n\) 矩阵，\(X\) 为 \(n \times s\) 矩阵。

\begin{itemize}
  \item[（1）] 证明：\(r(A)=n\) 当且仅当矩阵方程 \(A X=O\) 只有零解。
  \item[（2）] 若 \(r(A)<n, n-r(A) \leq s\) ，证明：\(A X=O\) 有解 \(X=B\) ，且满足 \(A \beta=0\) 的 \(\beta\) 都有 \(B X=\beta\) ．
\end{itemize}

\subsection{题目3}
3．设 \(C=\left(\begin{array}{cc}A & O \\ O & B\end{array}\right)\) ．

\begin{itemize}
  \item[（1）] 证明：\(C\) 可逆当且仅当 \(A, B\) 都可逆．\(C\) 可逆时，求 \(C\) 的逆．在一般情况下，求 \(C\) 的伴随．
  \item[（2）] 证明：\(C\) 正定当且仅当 \(A, B\) 都正定。
\end{itemize}

\subsection{题目4}
4．设 \(\varphi\) 是 \(U\) 到 \(V\) 的一个线性映射，且 \(\operatorname{dim} U=\operatorname{dim} V\) 。证明下列断言等价：

\begin{itemize}
  \item[（1）] \(\varphi\) 是满射．
  \item[（2）] \(\varphi\) 是单射．
  \item[（3）] 取 \(U\) 的一组基 \(\varepsilon_{1}, \varepsilon_{2}, \cdots, \varepsilon_{n}\) ，则 \(\varphi\left(\varepsilon_{1}\right), \varphi\left(\varepsilon_{2}\right), \cdots, \varphi\left(\varepsilon_{n}\right)\) 是 \(V\) 的一组基。
\end{itemize}

\subsection{题目5}
5．设 \(\mathcal{A}\) 是线性空间 \(V\) 上的线性变换， \(\mathcal{E}\) 是恒等变换．

\begin{itemize}
  \item[（1）] 证明： \(\mathcal{A}^{2}=\mathcal{E}\) 当且仅当 \(V=\operatorname{Ker}(\mathcal{E}-\mathcal{A}) \oplus \operatorname{Ker}(\mathcal{E}+\mathcal{A})\) 。
  \item[（2）] 存在正整数 \(k\) 使得 \(\mathcal{A}^{k}=\mathcal{E}\) ，证明：\(V=\operatorname{Ker}(\mathcal{E}-\mathcal{A}) \oplus \operatorname{Ker}\left(\mathcal{E}+\mathcal{A}+\cdots+\mathcal{A}^{k-1}\right)\) ．
\end{itemize}

\subsection{题目6}
6．设 \(A=E_{n}-\alpha \beta^{\mathrm{T}}\) ，其中 \(\alpha \beta^{\mathrm{T}} \neq O\) ，且 \(\beta^{\mathrm{T}} \alpha=c\) 。

\begin{itemize}
  \item[（1）] 计算 \(A^{2}\) ，并说明 \(A\) 不可逆等价于 \(c=1\) 。
  \item[（2）] 当 \(c \neq 1\) 时，求 \(A\) 的逆。
  \item[（3）] 当 \(c=0\) 时，求 \(A\) 的若尔当标准形，并针对不同的特征值 \(\lambda\) ，分别求对应的特征子空间．
\end{itemize}

\subsection{题目7}
7．在数域 \(\mathbb{K}\) 上的 \(n\) 阶矩阵空间中选取 \(A\) ，定义线性变换 \(\varphi(X)=A X-X A^{\mathrm{T}}\) 。

\begin{itemize}
  \item[（1）] 当 \(n=2, A=\left(\begin{array}{ll}0 & 1 \\ 0 & 0\end{array}\right)\) 时，求 \(\varphi\) 在 \(E_{11}, E_{12}, E_{21}, E_{22}\) 下的矩阵，并判断是否可对角化．
  \item[（2）] \(\alpha, \beta\) 为 \(A\) 的两个任意的特征向量，证明：\(\alpha \beta^{\mathrm{T}}\) 是 \(\varphi\) 的特征向量。
  \item[（3）] 若 \(A\) 有 \(n\) 个互异的特征值，证明：\(\varphi\) 可对角化，且 \(V=\operatorname{Ker}(\varphi) \oplus \operatorname{Im} \varphi\) ．
\end{itemize}

\subsection{题目8}
8．在次数小于 3 的多项式空间 \(V=\mathbb{R}[x]_{3}\) 中定义内积

\[
(f(x), g(x))=f(-1) g(-1)+f(0) g(0)+f(1) g(1) .
\]

\begin{itemize}
  \item[（1）] 求基 \(\alpha_{1}=1, \alpha_{2}=x, \alpha_{3}=x^{2}\) 的度量矩阵。
  \item[（2）] 求由 \(\alpha_{2}, \alpha_{3}\) 生成的子空间 \(L\left(\alpha_{2}, \alpha_{3}\right)\) 的正交补空间的基元 \(\beta_{1}\) 。
  \item[（3）] 求向量 \(\beta_{1}, \beta_{2}, \beta_{3}\) ，使得 \(\left(\alpha_{i}, \beta_{j}\right)=\delta_{i j}\) ，其中当 \(i=j\) 时，\(\delta_{i j}=1\) ，当 \(i \neq j\) 时，\(\delta_{i j}=0\) 。并证明\(\beta_{1}, \beta_{2}, \beta_{3}\) 也是 \(V\) 的一组基．
\end{itemize}
\end{document}
